\begin{textsample}{POS Dim 3 – 2023\_c5 – Score 4.00 – 2023\_c5/enade\_2023\_biomed\_7d177961fe.txt}  \label{ex:f3_pos_054}
CAPÍTULO 4 DISTRIBUIÇÃO DE CONCEITOS A finalidade deste capítulo é apresentar um panorama nacional sobre a distribuição do Conceito Enade 2023.
Para \textbf{traçar} esse panorama, recorre -se a tabelas e gráficos, que são acompanhados de comentários e análises que contribuem para a contextualização dos dados neles expostos.
Para visualizar o panorama do Enade 2023, em nível nacional, na primeira seção deste capítulo, os dados sobre o número de cursos participantes são organizados por grande região.
Na segunda, o mapeamento do conceito do exame é feito por categoria administrativa das IES e modalidade de oferta dos cursos, de maneira cruzada às grandes regiões.
Na última seção, a lógica de apresentação dos dados é feita por organização acadêmica e, também, por grande região de funcionamento dos cursos, no país como um todo. 4.1 PANORAMA NACIONAL DA DISTRIBUIÇÃO DOS CONCEITOS A lógica adotada na apresentação dos dados, nesta seção e nas duas seguintes, pressupõe o registro de duas informações de natureza reiterativa.
O primeiro diz respeito aos valores possíveis do Enade, que variam em níveis de 1 a 5, sendo que, à medida que esses valores aumentam, o curso é melhor avaliado, e ao fato de que aqueles cursos que tiveram somente um ou nenhum concluinte participante não obtêm o Conceito Enade, ficando sem conceito ( SC ), como ressaltado na apresentação deste relatório.
O segundo registro é que se houver a diferença entre os cursos tabulados neste capítulo e no capítulo 1 corresponderá aos cursos que não foram avaliados, ou seja, àqueles, em princípio, sem estudantes concluintes que tivessem sido inscritos para a prova.
Na Tabela 4.1, são apresentadas a quantidade e a distribuição de cursos de Biomedicina participantes do Enade 2023, por faixa de conceito e grande região.
A partir dos dados da Tabela 4.1, nota -se que, dos 503 cursos participantes, 181 ( 36,0 \% ) classificaram -se com conceito 2, o valor modal.
Esse foi o conceito modal nas seguintes regiões : Centro-Oeste ( 31,3 \% ), Norte ( 52,6 \% ) e Sudeste ( 39,2 \% ).
O conceito 3 foi o segundo mais frequente em nível nacional ( 31,6 \%, correspondendo a 159 cursos ), na maioria das regiões, exceto no Nordeste e Sul.
O conceito 4 foi o terceiro ( 12,7 \%, correspondendo a 64 cursos ).
Houve, ainda, 59 ( 11,7 \% ) cursos que receberam o conceito 1 e 25 ( 5,0 \% ) cursos que receberam o conceito 5.
Dos 503 cursos de Biomedicina, 15 ( 3,0 \% ) ficaram sem conceito ( SC ).
A região Norte participou com 38 cursos, o que corresponde a 7,6 \% do total nacional.
Desses, 20 cursos receberam o conceito 2 ( valor modal ), o que equivale a 52,6 \% do total regional, conforme mencionado.
A nove cursos ( 23,7 \% ) atribuiu -se o conceito 1 ; a sete cursos ( 18,4 \% ) o conceito 3 ; e a um curso ( 2,6 \% ) o conceito 5.
Ainda nessa região, um curso ( 2,6 \% ) obteve conceito 4 e nenhum curso ( 0,0 \% ) ficou sem conceito ( SC ).
A região Nordeste participou com 86 cursos, correspondentes a 17,1 \% do total nacional.
Nessa região, 31 cursos ( 36,0 \% em termos regionais ) obtiveram conceito 3, o conceito modal para a região.
Os conceitos 2 e 4 foram atribuídos, respectivamente, a 27 e a 11 cursos ( 31,4 \% e 12,8 \%, respectivamente ).
Ainda nessa região, sete cursos ( 8,1 \% ) foram avaliados com o conceito 1 e cinco cursos ( 5,8 \% ) foram avaliados com o conceito 5.
Outros cinco cursos ( 5,8 \% ) ficaram sem conceito ( SC ).
Para os 240 cursos participantes da região Sudeste ( 47,7 \% do total nacional ), o conceito 2 foi o conceito modal, tendo sido obtido por 94 cursos ( 39,2 \% ).
O conceito 3 foi atribuído a 75 cursos ( 31,3 \% ) e o conceito 4, a 29 cursos ( 12,1 \% ).
Receberam os conceitos 1 e 5, respectivamente, 22 cursos ( 9,2 \% ) e 13 cursos ( 5,4 \% ).
Nessa região, sete cursos ( 2,9 \% ) ficaram sem conceito ( SC ).
Os 91 cursos da região Sul corresponderam a 18,1 \% do total nacional.
Houve predominância do conceito 3 ( conceito modal ), atribuído a 35,2 \% dos cursos da região, ou seja, atribuído a 32 dos 91 cursos participantes na região Sul.
O conceito 2 foi atribuído a 25 cursos ( 27,5 \% ) ; o conceito 4, a 14 cursos ( 15,4 \% ) ; e o conceito 1, a 12 cursos ( 13,2 \% ).
Nessa região, seis cursos ( 6,6 \% ) obtiveram conceito 5 e dois cursos ( 2,2 \% ) ficaram sem conceito ( SC ).
Já dos 48 cursos participantes na região Centro-Oeste ( 9,5 \% do total nacional ), 15 ( 31,3 \% em termos regionais ) receberam o conceito 2, o conceito modal, como já destacado.
Os 14 cursos ( 29,2 \% ) obtiveram o conceito 3, nove cursos ( 18,8 \% ), o conceito 4 e nove cursos ( 18,8 \% ), o conceito 1.
Nessa região, nenhum curso ( 0,0 \% ) obteve o conceito 5 e um curso ( 2,1 \% ) ficou sem conceito ( SC ).
Em um gráfico de distribuição cumulativa, em cada valor de abscissa ( eixo x ), o valor apresentado no eixo das ordenadas ( eixo y ) é o valor acumulado da variável ( Conceito Enade ).
Em linhas gerais, em um gráfico desta mesma tipologia, poligonais mais à esquerda representam, no caso dos dados apresentados nesta seção, as grandes regiões com a pior distribuição, e poligonais mais à direita, as grandes regiões com os melhores conceitos.
A partir destas diretrizes, no Gráfico 4.1, é apresentada a distribuição cumulativa do Conceito Enade, segundo a grande região.
No Gráfico 4.1, observando a linha vertical pontilhada, que \textbf{ajuda} a visualizar o percentual de conceitos acumulados até o valor 2, ou seja, conceitos 1 e 2, pode -se notar, por exemplo, que a região Norte ( linha azul ) apresenta 76,3 \% dos cursos com conceito 2 ou abaixo deste valor.
A região Sul ( linha roxa ), por outro lado, apresenta 41,6 \% dos cursos na mesma situação : com conceito 2 ou abaixo.
Destaca -se que, na região Centro-Oeste ( linha \textbf{laranja} ), o conceito 5 não foi alcançado por nenhum curso, diferentemente das regiões Nordeste ( linha verde ), Norte ( linha azul ), Sudeste ( linha vermelha ) e Sul ( linha roxa ), nas quais todos os cursos têm conceito menor ou igual a 5, por isso as respectivas poligonais cumulativas dessas regiões alcançam 100 \% no conceito 5.
A região Norte ( linha azul ) apresenta os cursos com conceitos mais baixos, ao passo que a região Sul ( linha roxa ) apresenta a distribuição com valores maiores.
Os cursos em IES situadas na região Sudeste ( linha vermelha ) apresentam uma situação intermediária. 4.2 CONCEITOS POR CATEGORIA ADMINISTRATIVA, MODALIDADE DE OFERTA DOS CURSOS E GRANDE REGIÃO Na Tabela 4.2, é apresentada a distribuição dos cursos de Biomedicina participantes do Enade 2023, por categoria administrativa e modalidade de oferta, de acordo com os conceitos por eles alcançados, segundo as grandes regiões brasileiras.
De acordo com as informações da Tabela 4.2, dos 503 cursos participantes, 31 ( 6,2 \% ) eram ministrados em instituições públicas e 472 ( 93,8 \% ), em instituições privadas.
Quanto à modalidade de oferta, 456 ( 90,7 \% ) cursos eram presenciais e 47 ( 9,3 \% ), a distância.
De acordo com as informações da mesma tabela, em termos nacionais, entre instituições públicas e privadas, observa -se que, dos 25 cursos avaliados com conceito 5, 14 foram oferecidos em IES públicas, e 11, em IES privadas.
Dos 472 cursos participantes de IES privadas, o conceito 2 foi o valor modal, atribuído a 180 ( 38,1 \% ) cursos.
Em relação aos demais cursos de instituições privadas, 155 cursos ( 32,8 \% ) foram avaliados com o conceito 3, 59 cursos ( 12,5 \% ) foram avaliados com o conceito 1, 52 cursos ( 11,0 \% ) foram avaliados com o conceito 4 e, 11 cursos ( 2,3 \% ), com o conceito 5.
Nessa categoria, 15 cursos ( 3,2 \% ) ficaram sem conceito ( SC ).
Na rede pública, o conceito modal foi o 5, atribuído a 14 cursos, correspondentes a 45,2 \% dos 31 cursos da categoria.
Dos demais cursos participantes, 12 cursos ( 38,7 \% ) obtiveram o conceito 4, quatro cursos ( 12,9 \% ) receberam o conceito 3 e um curso ( 3,2 \% ), o conceito 2.
O conceito 1 não foi atribuído a nenhum curso ( 0,0 \% ), conforme mencionado.
Nessa categoria administrativa, nenhum curso ( 0,0 \% ) ficou sem conceito ( SC ).
Dos 456 cursos participantes da modalidade presencial, o conceito 2 foi o valor modal, atribuído a 160 ( 35,1 \% ) cursos.
Dos demais cursos dessa modalidade de oferta, 155 cursos ( 34,0 \% ) receberam o conceito 3 ; 64 cursos ( 14,0 \% ) receberam o conceito 4 ; 40 cursos ( 8,8 \% ), o conceito 1 ; e 25 cursos ( 5,5 \% ), o conceito 5.
Nessa modalidade de oferta, 12 cursos ( 2,6 \% ) ficaram sem conceito ( SC ).
Dos 47 cursos participantes da modalidade a distância, o conceito 2 foi o valor modal, atribuído a 21 ( 44,7 \% ) cursos.
Dos demais cursos ofertados nessa modalidade, 19 cursos ( 40,4 \% ) receberam o conceito 1, quatro cursos ( 8,5 \% ) receberam o conceito 3 ; nenhum curso ( 0,0 \% ) recebeu o conceito 5 ; e nenhum curso ( 0,0 \% ) recebeu o conceito 4.
Nessa modalidade de oferta, três cursos ( 6,4 \% ) ficaram sem conceito ( SC ).
Na análise por região, observa -se, que, na região Norte, as instituições públicas participaram com um curso ( 2,6 \% total regional ).
O conceito modal foi 4, atribuído a um curso, correspondendo a 100,0 \% dessa categoria na região.
Os demais conceitos não foram atribuídos a nenhum curso.
As instituições privadas participaram com 37 cursos na região Norte ( 97,4 \% do total regional ), e nenhum curso ficou sem conceito ( SC ).
O conceito modal nas instituições privadas na região foi o 2, atribuído a 20 cursos, o que corresponde a 54,1 \% dessa categoria na região.
Os conceitos 1, 3 e 5 foram atribuídos a nove, sete e um cursos, respectivamente.
O conceito 4 não foi atribuído a nenhum curso ( 0,0 \% ).
Nenhum curso ficou sem conceito ( SC ) ou recebeu os demais conceitos.
Dos cursos oferecidos na região Norte, 35 eram presenciais e nenhum ficou sem conceito ( SC ).
Dos 35 cursos oferecidos, 19 obtiveram o conceito 2 ( modal ) ; sete, o conceito 3 ; sete, o conceito 1 ; um, o conceito 5 ; e um, o conceito 4.
Nenhum curso ficou sem conceito ( SC ) ou recebeu os demais conceitos na região Norte e na modalidade presencial.
Na modalidade a distância, havia três cursos, tendo dois recebido o conceito modal 1 ; um, o conceito 2 ; nenhum, o conceito 5 ; nenhum, o conceito 4 ; e nenhum, o conceito 3.
Nenhum curso ficou sem conceito ( SC ) ou recebeu os demais conceitos na região Norte e na modalidade a distância.
Na região Nordeste, a rede privada concentra 81 dos 86 cursos participantes, o equivalente a 94,2 \% do total da região.
Entre os cursos oferecidos por IES privadas no Nordeste, 31 cursos foram avaliados com o conceito 3, o conceito modal ; 27 cursos obtiveram o conceito 2 ; nove cursos, o conceito 4 ; sete cursos, o conceito 1 ; e dois cursos, o conceito 5.
Nessa categoria administrativa e grande região, cinco cursos ficaram sem conceito ( SC ).
As instituições públicas da região Nordeste participaram com cinco cursos ( 5,8 \% do total da região ).
Desses, três ( 60,0 \% ) obtiveram o conceito 5, o valor modal, seguidos por dois cursos, que receberam o conceito 4.
Os conceitos 3, 2 e 1 não foram atribuídos a nenhum curso.
Nenhum curso ficou sem conceito ( SC ).
Dos cursos oferecidos na região Nordeste, 83 eram presenciais : cinco ficaram sem conceito ( SC ), quatro obtiveram o conceito 1 ; 27, o conceito 2 ; 31, o conceito 3, o valor modal ; 11, o conceito 4 ; e cinco, o conceito 5.
Na modalidade a distância, eram três cursos, tendo os três recebido o conceito modal 1.
Os conceitos 2, 3, 4 e 5 não foram atribuídos a nenhum curso.
Nenhum curso ficou sem conceito ( SC ).
Na região Sudeste, a proporção de cursos da rede privada, 93,8 \%, correspondeu a 225 dos 240 cursos participantes.
Nessa combinação de categoria administrativa e grande região, o conceito modal foi o 2, atribuído a 93 cursos ; e sete cursos ficaram sem conceito ( SC ).
Os demais foram avaliados com os seguintes conceitos : 3 ( 73 cursos ), 4 ( 23 cursos ), 1 ( 22 cursos ) e 5 ( sete cursos ).
Dos 15 cursos oferecidos por instituições públicas na região Sudeste, seis obtiveram o conceito 5 ; seis, o conceito 4 ; dois, o conceito 3 ; um obteve o conceito 2 e nenhum obteve o conceito 1.
Nenhum curso ficou sem conceito ( SC ).
Ainda na região Sudeste, grande parcela dos cursos ( 91,3 \% ) foi oferecida na modalidade presencial : 219 dos 240 cursos.
Para esses, a moda foi o conceito 2, com 82 cursos.
Os demais cursos presenciais da região foram alocados nos conceitos 3 ( 73 cursos ), 4 ( 29 cursos ), 1 ( 16 cursos ) e 5 ( 13 cursos ).
Ficaram sem conceito ( SC ) seis cursos presenciais.
Na região Sudeste, 21 cursos foram oferecidos na modalidade a distância, tendo sido seis com o conceito 1 ; 12 com o conceito 2 ; dois com o conceito 3 ; nenhum com o conceito 4 ; e nenhum com o conceito 5.
Ficou sem conceito ( SC ) um curso.
As instituições privadas concentraram 84 dos 91 cursos participantes da região Sul, 92,3 \% do total regional.
Desses, 31 cursos obtiveram o conceito 3, o conceito modal.
Nessa combinação de categoria administrativa e grande região, 25 receberam o conceito 2 ; 13 o conceito 4 ; 12, o conceito 1 ; e um obteve o conceito 5.
Dois cursos ficaram sem conceito ( SC ).
As instituições públicas na região Sul participaram com sete cursos ( 7,7 \% dos cursos da região ), aos quais foram atribuídos os conceitos 1 ( nenhum curso ), 2 ( nenhum curso ), 3 ( um curso ), 4 ( um curso ) e 5 ( cinco cursos ).
Nenhum curso da região Sul oferecido por IES públicas ficou sem conceito ( SC ).
Na região Sul, 72 dos cursos foram oferecidos na modalidade presencial.
Para esses, a moda foi o conceito 3, com 31 cursos.
Os demais cursos presenciais da região foram alocados nos conceitos 2 ( 17 cursos ), 4 ( 14 cursos ), 5 ( seis cursos ) e 1 ( quatro cursos ).
Ficaram sem conceito ( SC ) nenhum curso.
Na modalidade a distância, foram oferecidos 19 cursos, sendo oito cursos com conceito 1, oito cursos com conceito 2, um curso com conceito 3, nenhum curso com conceito 4 e nenhum curso com conceito 5.
Dois cursos ficaram sem conceito ( SC ) nessa região.
Na região Centro-Oeste, 45 dos 48 cursos participantes eram de IES privadas ( 93,8 \% em termos regionais ).
Desses, 15 cursos receberam conceito 2, o valor modal.
Dos demais cursos, 13 receberam o conceito 3, nove, o conceito 1, sete, o conceito 4 e nenhum, o conceito 5.
Um curso ficou sem conceito ( SC ).
Quanto aos três cursos oferecidos por instituições públicas na região Centro-Oeste ( 6,3 \% do total regional ), nenhum ficou com conceito 1 ; nenhum, com conceito 2 ; um, com conceito 3 ; dois, com conceito 4 ; e nenhum ficou com conceito 5.
Nenhum curso ficou sem conceito ( SC ) nessa região.
Também na região Centro-Oeste, a maioria dos cursos foi ofertada na modalidade presencial : 47 ( 97,9 \% ).
Para esses, a moda foi o conceito 2, com 15 cursos ( 31,9 \% ).
Os demais cursos presenciais da região foram alocados nos conceitos 3 ( 13 cursos, 27,7 \% ), 4 ( nove cursos, 19,1 \% ), 1 ( nove cursos, 19,1 \% ) e 5 ( nenhum curso, 0,0 \% ).
Um curso presencial da região ficou sem conceito ( SC ).
Na modalidade a distância, foi oferecido um curso, sendo nenhum curso com conceito 1, nenhum curso com conceito 2, um curso com conceito 3 ( modal ), nenhum curso com conceito 4 e nenhum curso com conceito 5.
Nenhum curso ficou sem conceito ( SC ) nessa região.
A título de síntese, no Gráfico 4.2, apresenta -se a distribuição do Conceito Enade, segundo a categoria administrativa da IES.
Nesse gráfico, como dito anteriormente, poligonais mais à esquerda representam Conceito Enade mais baixos, e poligonais mais à direita, Conceito Enade mais altos.
É importante observar que, para cada poligonal presente nesse gráfico, a porcentagem constante no eixo y refere -se ao universo de cursos oferecidos por instituições públicas ( poligonal azul ) ou privadas ( poligonal verde ), respectivamente.
Os cursos em IES públicas ( poligonal azul ) apresentam uma poligonal mais à direita do que os em IES privadas ( poligonal verde ) e, consequentemente, uma distribuição de Conceito Enade com valores maiores. 4.3 CONCEITOS POR ORGANIZAÇÃO ACADÊMICA E GRANDE REGIÃO Na Tabela 4.3, encontra -se a distribuição dos conceitos atribuídos aos cursos participantes do Enade 2023, na área de Biomedicina, por organização acadêmica, segundo as grandes regiões brasileiras.
Como mostram os dados da Tabela 4.3, dos 503 cursos de Biomedicina participantes, 170 eram oferecidos em universidades ; 209, em centros universitários ; 123, em faculdades ; e um curso, por CEFET/IF.
Essa distribuição corresponde a, respectivamente, 33,8 \%, 41,6 \%, 24,5 \% e 0,2 \% dos cursos.
De acordo com os dados apresentados na mesma tabela, de todos os 181 cursos avaliados com o conceito 2, 43 eram vinculados a universidades.
Os cursos vinculados a esse tipo de organização acadêmica tiveram o conceito 3 como conceito modal, atribuído a 58 cursos ( 34,1 \% ).
Os demais cursos ofertados por universidades receberam os conceitos 4 ( 31 cursos ), 5 ( 19 cursos ) e 1 ( 16 cursos ).
Três cursos ficaram sem conceito ( SC ).
Dos cursos vinculados a centros universitários, o conceito modal foi o 2, atribuído a 84 cursos ( 40,2 \% ).
Oito cursos vinculados a esse tipo de organização acadêmica ficaram sem conceito ( SC ), e os outros cursos receberam os conceitos 3 ( 67 cursos ), 4 ( 26 cursos ), 1 ( 21 cursos ) e 5 ( três cursos ).
Dos 123 cursos mantidos por faculdades, 54 ( 43,9 \% ) receberam o conceito 2, o conceito modal.
O segundo conceito mais frequentemente atribuído aos cursos oferecidos por essa categoria de IES foi o conceito 3, recebido por 34 cursos.
Na sequência, 22 cursos obtiveram o conceito 1 ; seis cursos, o conceito 4 ; e três cursos, o conceito 5.
Dos cursos oferecidos em faculdades, quatro cursos ficaram sem conceito ( SC ).
Dos cursos oferecidos por CEFET/IF, o conceito modal foi o 4, atribuído a um curso ( 100,0 \% ).
Nenhum curso vinculado a esse tipo de organização acadêmica ficou sem conceito ( SC ), e os outros cursos receberam os conceitos 5 ( nenhum curso ), 3 ( nenhum curso ), 2 ( nenhum curso ) e 1 ( nenhum curso ).
Considerando -se, separadamente, as regiões brasileiras, verifica -se que, na região Norte, as universidades participaram com sete cursos avaliados, aos quais foram atribuídos os seguintes conceitos : 1 a dois cursos ; 2 a dois cursos ; 3 a dois cursos ; 4 a um curso ; e 5 a nenhum curso.
Nenhum curso ficou sem conceito ( SC ).
Os centros universitários da região Norte foram representados por 18 cursos, dos quais três cursos obtiveram o conceito 1 ; nove cursos, o conceito 2 ; cinco cursos, o conceito 3 ; nenhum curso, o conceito 4 ; e um curso, o conceito 5.
Nenhum curso ficou sem conceito ( SC ).
As faculdades participaram com 13 cursos na região Norte, os quais receberam os conceitos 1 ( quatro cursos ), 2 ( nove cursos ), 3 ( nenhum curso ), 4 ( nenhum curso ) e 5 ( nenhum curso ).
Nenhum curso ficou sem conceito ( SC ).
Os CEFET/IF participaram com nenhum curso, aos quais foram atribuídos os seguintes conceitos : 1 a nenhum curso ; 2 a nenhum curso ; 3 a nenhum curso ; 4 a nenhum curso ; e 5 a nenhum curso.
Nenhum curso ficou sem conceito ( SC ).
Na região Nordeste, as universidades participaram com 15 dos 86 cursos da área de Biomedicina oferecidos.
O conceito modal foi o 3, atribuído a quatro cursos.
Outros três cursos obtiveram o conceito 5 ; três cursos, o conceito 4 ; três cursos, o conceito 2 ; e dois cursos, o conceito 1.
Nenhum curso ficou sem conceito ( SC ).
Os centros universitários contaram com 49 cursos participantes na região Nordeste, 19 dos quais receberam o conceito modal 3.
Para os demais cursos, foram atribuídos os conceitos 2 ( 14 cursos ), 4 ( sete cursos ), 1 ( três cursos ) e 5 ( dois cursos ).
Quatro cursos oferecidos por centros universitários no Nordeste ficaram sem conceito ( SC ).
As faculdades foram representadas por 22 cursos na região Nordeste, os quais receberam os conceitos 1 ( dois cursos ), 2 ( dez cursos ), 3 ( oito cursos ), 4 ( um curso ) e 5 ( nenhum curso ).
Um curso ficou sem conceito ( SC ).
Na região Nordeste, os CEFET/IF participaram com nenhum curso, aos quais foram atribuídos os seguintes conceitos : 1 a nenhum curso ; 2 a nenhum curso ; 3 a nenhum curso ; 4 a nenhum curso ; e 5 a nenhum curso.
Nenhum curso ficou sem conceito ( SC ).
Na região Sudeste, as universidades concentraram 104 dos 240 cursos de Biomedicina da região.
Dos cursos oferecidos em IES com este tipo de organização acadêmica, nessa região, o conceito modal foi 2, atribuído a 33 cursos.
Os demais cursos receberam os conceitos 3 ( 32 cursos ), 4 ( 17 cursos ), 5 ( dez cursos ) e 1 ( nove cursos ).
Três cursos ficaram sem conceito ( SC ).
Os centros universitários participaram com 83 cursos na região Sudeste, dos quais 37 obtiveram o conceito modal 2.
O conceito 3 foi atribuído a 27 cursos ; o conceito 4 a nove cursos ; o conceito 1 a oito cursos ; e o conceito 5 a nenhum curso.
Dois cursos ficaram sem conceito ( SC ) nessa combinação de organização acadêmica e região.
Já as faculdades foram representadas por 52 cursos na região Sudeste, dos quais dois ficaram sem conceito ( SC ).
O conceito modal 2 foi atribuído a 24 cursos.
Aos demais cursos foram atribuídos os conceitos 3 ( 16 cursos ), 1 ( cinco cursos ), 5 ( três cursos ) e 4 ( dois cursos ).
Os CEFET/IF participaram com um curso na região Sudeste, aos quais foram atribuídos os seguintes conceitos : 1 a nenhum curso ; 2 a nenhum curso ; 3 a nenhum curso ; 4 a um curso ; e 5 a nenhum curso.
Nenhum curso ficou sem conceito ( SC ).
Dos 91 cursos da região Sul, 33 eram mantidos por universidades, para os quais o conceito modal foi 3, atribuído a 15 cursos.
Os demais cursos receberam os conceitos 5 ( seis cursos ), 4 ( seis cursos ), 2 ( três cursos ) e 1 ( três cursos ).
Em IES com esse tipo de organização acadêmica, nenhum curso da região Sul ficou sem conceito ( SC ).
Dos 42 cursos participantes de centros universitários da região Sul, dois ficaram sem conceito ( SC ), e 17 cursos receberam conceito 2, o valor modal.
Os conceitos restantes ficaram assim distribuídos : 3 ( 11 cursos ), 4 ( seis cursos ), 1 ( seis cursos ) e 5 ( nenhum curso ).
Quanto aos 16 cursos vinculados a faculdades na região Sul, seis cursos receberam o conceito 3, valor modal.
Os demais conceitos ficaram assim distribuídos : 2 ( cinco cursos ), 1 ( três cursos ), 4 ( dois cursos ) e 5 ( nenhum curso ).
Dos 16 cursos, nenhum ficou sem conceito ( SC ).
Na região Centro-Oeste, 11 dos 48 cursos oferecidos eram mantidos por universidades.
Para esse tipo de organização acadêmica, o conceito modal foi 3, atribuído a cinco cursos.
Em relação às universidades, nenhum curso da região ficou sem conceito ( SC ).
Os demais cursos receberam os conceitos 4 ( quatro cursos ), 2 ( dois cursos ), 5 ( nenhum curso ) e 1 ( nenhum curso ).
Os centros universitários da região Centro-Oeste contaram com 17 cursos : sete cursos foram avaliados com o conceito 2 ( modal ) ; cinco cursos, com o conceito 3 ; quatro cursos, com o conceito 4 ; um curso, com o conceito 1 ; e nenhum curso obteve o conceito 5.
Nenhum curso ficou sem conceito ( SC ).
Dos 20 cursos oferecidos em faculdades na região Centro-Oeste, oito cursos receberam o conceito 1 ( modal ).
Os conceitos restantes ficaram assim distribuídos : 2 ( seis cursos ) ; 3 ( quatro cursos ) ; 4 ( um curso ) ; e 5 ( nenhum curso ).
Nessa região, um curso oferecido em faculdades ficou sem conceito ( SC ).
No que se refere aos CEFET/IF da região Centro-Oeste, nenhum curso foi avaliado : nenhum curso deles com o conceito 5 ; nenhum curso com o conceito 4 ; nenhum curso com o conceito 3 ; nenhum curso com o conceito 2 ; nenhum curso com o conceito 1 ; e nenhum curso ficou sem conceito ( SC ).
De maneira semelhante à forma como se procedeu em relação à seção anterior, esta seção \textbf{encerra} -se com o Gráfico 4.3, que apresenta a distribuição do Conceito Enade dos cursos de Biomedicina segundo a organização acadêmica da IES.
No Gráfico 4.3, apresenta -se a distribuição do Conceito Enade segundo a organização acadêmica da IES.
A leitura desse último gráfico é semelhante à dos anteriores, seguindo -se a lógica de que poligonais deslocadas para a esquerda correspondem aos conceitos mais baixos ( 1, 2 e 3 ) e poligonais deslocadas para a direita correspondem a Conceitos Enade mais altos.
Nas universidades ( linha azul ), todos os cursos obtiveram conceito 5 ou abaixo disso, já que a poligonal cumulativa alcança 100 \% para o conceito 5.
Além disso, os cursos de universidades apresentam uma poligonal mais à direita do que as demais, denotando uma melhor distribuição de Conceitos Enade.
No outro extremo, encontram -se os cursos mantidos pelas faculdades ( linha vermelha ), apresentando as respectivas poligonais mais à esquerda e, consequentemente, uma distribuição de Conceitos Enade com os menores valores.
Os cursos mantidos por centros universitários ( linha verde ) apresentam poligonais mais centrais e, consequentemente, uma distribuição de Conceitos Enade com os valores intermediários.

% matched lemmas: ajudar, encerrar, laranja, traçar
\end{textsample}
